%-------------------------
% classic-latex: one-page, single-column resume.
% Body is filled from the resume document by engine/render.py:
%   see engine/render.py for the template syntax (VAR blocks and %%-statements).
%------------------------
\documentclass[letterpaper,11pt]{article}
\usepackage{latexsym}
\usepackage[empty]{fullpage}
\usepackage{titlesec}
\usepackage{marvosym}
\usepackage[usenames,dvipsnames]{color}
\usepackage{verbatim}
\usepackage{enumitem}
\usepackage[hidelinks]{hyperref}
\usepackage{fancyhdr}
\usepackage[english]{babel}
\usepackage{tabularx}
\usepackage{fontawesome5}
\usepackage{multicol}
\setlength{\multicolsep}{-3.0pt}
\setlength{\columnsep}{-1pt}
\input{glyphtounicode}
\pagestyle{fancy}
\fancyhf{}
\fancyfoot{}
\renewcommand{\headrulewidth}{0pt}
\renewcommand{\footrulewidth}{0pt}
\addtolength{\oddsidemargin}{-0.6in}
\addtolength{\evensidemargin}{-0.5in}
\addtolength{\textwidth}{1.19in}
\addtolength{\topmargin}{-.7in}
\addtolength{\textheight}{1.4in}
\urlstyle{same}
\raggedbottom
\raggedright
\setlength{\tabcolsep}{0in}
\titleformat{\section}{
  \vspace{-6pt}\scshape\raggedright\large\bfseries
}{}{0em}{}[\color{black}\titlerule \vspace{-5pt}]
\pdfgentounicode=1
\newcommand{\resumeItem}[1]{\small\item{{#1 \vspace{-2pt}}}}
\newcommand{\resumeSubheading}[4]{
  \item
    \begin{tabular*}{1.0\textwidth}[t]{l@{\extracolsep{\fill}}r}
      \textbf{#1} & \textbf{\small #2} \\
      \textit{\small#3} & \textit{\small #4} \\
    \end{tabular*}\vspace{-6pt}
}
\newcommand{\resumeProjectHeading}[2]{
    \item
    \begin{tabular*}{1.001\textwidth}{l@{\extracolsep{\fill}}r}
      \small#1 & \textbf{\small #2}\\
    \end{tabular*}\vspace{-4pt}
}
\newcommand{\resumeProjectSubheading}[1]{
    \vspace{4pt}
    \item[]
    \small{\textbf{#1}}
    \vspace{-4pt}
}
\newcommand{\resumeItemListStart}{\begin{itemize}}
\newcommand{\resumeItemListEnd}{\end{itemize}\vspace{-5pt}}
\newcommand{\resumeSubHeadingListStart}{\begin{itemize}[leftmargin=0.0in, label={}]}
\newcommand{\resumeSubHeadingListEnd}{\end{itemize}}
\newcommand{\resumeSkillItem}[2]{
  \resumeItem{
    \textbf{#1}{:\hspace{0.5mm}#2 \vspace{-0.5mm}}
  }
}
\newcommand{\resumeSubItem}[2]{\resumeSkillItem{#1}{#2}\vspace{-4pt}}
\begin{document}
\begin{center}
    {\Huge \scshape Meera Rao} \\ \vspace{6pt}
    \small +91-90000-00000 \textbar{}
    \href{mailto:meera.rao@example.com}{\underline{meera.rao@example.com}} \textbar{}
    \href{https://www.linkedin.com/in/meera-rao-example/}{\underline{linkedin.com/in/meera-rao-example}} \textbar{}
    \href{https://meera-rao.example.com}{\underline{meera-rao.example.com (portfolio)}}
\end{center}
\section{Summary}
Mechanical design engineer with 5 years in automotive brake components: SolidWorks design of sheet-metal and cast parts, GD\&T drawings, tolerance stack-ups and ANSYS FEA, from DFMEA through engineering release in Teamcenter.
\section{Skills}
\resumeItemListStart[label=\textbullet]
  \resumeSubItem{CAD}{ SolidWorks, AutoCAD, Teamcenter, CATIA V5 }
  \resumeSubItem{Design \& analysis}{ GD\&T, Tolerance stack-up, ANSYS Mechanical, FEA, Sheet metal design, Casting design }
  \resumeSubItem{Quality \& process}{ DFMEA, DFM/DFA, Six Sigma }
\resumeItemListEnd
\section{Experience}
\resumeSubHeadingListStart
  \resumeSubheading{Design Engineer}{Pune}{Kalyani Precision Components}{Aug 2022 -- Present}
    \resumeItemListStart[label=\textbullet]
      \resumeItem{Own sheet-metal and casting design of brake-system brackets in SolidWorks, releasing GD\&T drawings (ASME Y14.5) and engineering changes through Teamcenter}
      \resumeItem{Redesigned a mounting bracket using ANSYS Mechanical FEA, cutting mass by 18\% while passing the same load cases}
      \resumeItem{Owned DFMEA for 2 new product lines, running design reviews with quality and manufacturing teams}
      \resumeItem{Trained 3 interns on drawing standards and GD\&T}
    \resumeItemListEnd
  \resumeSubheading{Graduate Engineer Trainee}{Pune}{Kalyani Precision Components}{Aug 2021 -- July 2022}
    \resumeItemListStart[label=\textbullet]
      \resumeItem{Rotated through manufacturing, quality and design; designed 5 drilling jigs and fixtures for the machining line}
    \resumeItemListEnd
\resumeSubHeadingListEnd
\section{Education}
\resumeSubHeadingListStart
  \resumeSubheading{B.E.\ in Mechanical Engineering}{Graduated: 2021}{Savitribai Phule Pune University}{CGPA 8.1 / 10}
\resumeSubHeadingListEnd
\section{Certifications}
\resumeItemListStart[label=\textbullet]
  \resumeItem{Six Sigma Green Belt (2023)}
  \resumeItem{Certified SolidWorks Associate (CSWA, 2020)}
\resumeItemListEnd
\end{document}
